% Generado por E2/informe/generar_anexo_c.py desde E2/src/lib/esquema.py. No editar a mano.

\section*{Anexo C. Diccionario de datos}
\addcontentsline{toc}{section}{Anexo C. Diccionario de datos}

Sale del mismo esquema que usa la cadena para construir el oro, así que no puede quedar desfasado del código. Para cada campo: su rol en el modelo, si se puede sumar, su unidad, su significado de negocio y de qué columna de origen sale y con qué regla. Las columnas nuevas de la E2 son las del dólar a la cobranza y al embarque, la marca de costo imputado y las cajas y el monto tal como venían en el ERP (sección~\ref{sec:modelo}).

\subsubsection*{\texttt{fact\_venta}}
Hecho transaccional. Una fila por línea de venta facturada.

{\footnotesize
\begin{longtable}{@{}>{\raggedright\arraybackslash}p{0.18\textwidth}>{\raggedright\arraybackslash}p{0.09\textwidth}>{\raggedright\arraybackslash}p{0.08\textwidth}>{\raggedright\arraybackslash}p{0.075\textwidth}>{\raggedright\arraybackslash}p{0.2\textwidth}>{\raggedright\arraybackslash}p{0.245\textwidth}@{}}
\toprule
Columna & Rol & Aditiv. & Unidad & Descripción & Origen y regla \\
\midrule
\endhead
\texttt{venta\_id} & degenerada & — & — & Número de la línea en el ERP. & ventas.\allowbreak{}venta\_id; deduplicado tras normalizar formatos \\
\texttt{fecha\_key} & llave & — & — & Fecha de factura. & ventas.\allowbreak{}fecha, leída en sus 4 formatos \\
\texttt{producto\_anada\_key} & llave & — & — & Producto y añada vendidos. & ventas.\allowbreak{}producto\_id + ventas.\allowbreak{}anada \\
\texttt{cliente\_key} & llave & — & — & Cliente que compra. & ventas.\allowbreak{}cliente\_id \\
\texttt{mercado\_key} & llave & — & — & Mercado de destino. & ventas.\allowbreak{}mercado\_id (coincide con clientes.\allowbreak{}mercado\_id) \\
\texttt{cajas} & medida & aditiva & cajas & Cajas vendidas; negativas si la línea es merma. & ventas.\allowbreak{}cajas; si vienen en botellas (12 o 10 veces las cajas del embarque), se dividen por el factor \\
\texttt{litros} & medida & aditiva & L & Litros vendidos. & cajas $\times$ dim\_producto.\allowbreak{}litros\_por\_caja \\
\texttt{monto\_bruto\_clp} & medida & aditiva & CLP & Monto antes de descuento, en pesos. & cajas $\times$ precio unitario; si la moneda es USD, $\times$ tipo\_cambio \\
\texttt{monto\_bruto\_usd} & medida & aditiva & USD & Monto antes de descuento, en dólares. & cajas $\times$ precio unitario; si la moneda es CLP, $\div$ tipo\_cambio \\
\texttt{monto\_neto\_clp} & medida & aditiva & CLP & Ingreso neto de descuento, en pesos. & cajas $\times$ precio unitario $\times$ (1 $-$ descuento), convertido con tipo\_cambio \\
\texttt{monto\_neto\_usd} & medida & aditiva & USD & Ingreso neto de descuento, en dólares. & cajas $\times$ precio unitario $\times$ (1 $-$ descuento), convertido con tipo\_cambio \\
\texttt{costo\_total\_clp} & medida & aditiva & CLP & Costo de producción de lo vendido, en pesos. & \textbar{}cajas\textbar{} $\times$ dim\_producto.\allowbreak{}costo\_caja\_clp; en la merma el costo no se reversa \\
\texttt{tipo\_cambio} & contexto & no aditiva & CLP/USD & Dólar observado usado para convertir la línea. & API mindicador.\allowbreak{}cl, fecha de factura; fines de semana y feriados con el último día hábil previo \\
\texttt{es\_merma} & indicador & aditiva & 0/1 & 1 si la línea es merma (cajas negativas). & cajas < 0; merma según la contraparte (4 sep) \\
\texttt{monto\_recalculado} & indicador & aditiva & 0/1 & 1 si el monto registrado no cuadraba y se recalculó. & \textbar{}monto registrado $-$ calculado\textbar{} > 0,5 \% \\
\texttt{cajas\_corregidas} & indicador & aditiva & cajas & 1 si las cajas venían en botellas. & cruce con origen\_embarques.\allowbreak{}json \\
\texttt{moneda\_imputada} & indicador & aditiva & 0/1 & 1 si la moneda venía vacía. & ventas.\allowbreak{}moneda\_id nulo; se imputa por mercado \\
\texttt{costo\_imputado} & indicador & aditiva & 0/1 & 1 si el costo de la caja se imputó por tier y añada (no venía en el maestro). & origen\_costos\_productos.\allowbreak{}csv sin costo para ese producto y añada \\
\texttt{cajas\_registradas} & medida & aditiva & cajas & Cajas tal como vienen en el ERP, antes de corregir las cargadas en botellas. & ventas.\allowbreak{}cajas sin corregir; solo para el escenario sin correcciones \\
\texttt{monto\_neto\_registrado\_usd} & medida & aditiva & USD & Monto neto tal como viene en el ERP, en dólares, antes de recalcularlo. & ventas.\allowbreak{}monto\_neto sin recalcular; solo para el escenario sin correcciones \\
\texttt{tipo\_cambio\_cobranza} & contexto & no aditiva & CLP/USD & Dólar observado a la fecha estimada de cobranza. & fecha de factura + 60 días (hemisferio sur) o + 90 (norte); misma regla de días hábiles \\
\texttt{tipo\_cambio\_embarque} & contexto & no aditiva & CLP/USD & Dólar observado a la fecha de embarque (fecha poco confiable). & origen\_embarques.\allowbreak{}json; solo para la sensibilidad, no para la métrica \\
\bottomrule
\end{longtable}
}

\subsubsection*{\texttt{dim\_tiempo}}
Una fila por día calendario de 2025.

{\footnotesize
\begin{longtable}{@{}>{\raggedright\arraybackslash}p{0.18\textwidth}>{\raggedright\arraybackslash}p{0.09\textwidth}>{\raggedright\arraybackslash}p{0.08\textwidth}>{\raggedright\arraybackslash}p{0.075\textwidth}>{\raggedright\arraybackslash}p{0.2\textwidth}>{\raggedright\arraybackslash}p{0.245\textwidth}@{}}
\toprule
Columna & Rol & Aditiv. & Unidad & Descripción & Origen y regla \\
\midrule
\endhead
\texttt{fecha\_key} & llave & — & — & Llave AAAAMMDD. & generada \\
\texttt{fecha} & atributo & — & fecha & Fecha. & calendario \\
\texttt{dia\_semana} & atributo & — & — & Lunes a domingo. & calendario \\
\texttt{semana} & atributo & — & — & Semana del año. & calendario \\
\texttt{mes} & atributo & — & — & Mes (número y nombre). & calendario \\
\texttt{trimestre} & atributo & — & — & Trimestre. & calendario \\
\texttt{semestre} & atributo & — & — & Primer o segundo semestre. & calendario \\
\texttt{anio} & atributo & — & año & Año. & calendario \\
\texttt{es\_dia\_habil} & atributo & — & 0/1 & 1 si es día hábil. & 1 si la API del dólar publica valor ese día \\
\texttt{es\_feriado} & atributo & — & 0/1 & 1 si es feriado. & día de semana sin valor publicado en la API \\
\bottomrule
\end{longtable}
}

\subsubsection*{\texttt{dim\_producto}}
Una fila por producto y añada (24 productos $\times$ 4 añadas = 96 filas).

{\footnotesize
\begin{longtable}{@{}>{\raggedright\arraybackslash}p{0.18\textwidth}>{\raggedright\arraybackslash}p{0.09\textwidth}>{\raggedright\arraybackslash}p{0.08\textwidth}>{\raggedright\arraybackslash}p{0.075\textwidth}>{\raggedright\arraybackslash}p{0.2\textwidth}>{\raggedright\arraybackslash}p{0.245\textwidth}@{}}
\toprule
Columna & Rol & Aditiv. & Unidad & Descripción & Origen y regla \\
\midrule
\endhead
\texttt{producto\_anada\_key} & llave & — & — & Llave producto-añada. & generada \\
\texttt{producto} & atributo & — & — & Nombre del vino. & origen\_costos\_productos.\allowbreak{}csv: producto \\
\texttt{cepa} & atributo & — & — & Cepa, con el nombre oficial del catálogo del ERP. & CSV: cepa, normalizada desde sus alias y conformada a origen\_vina.\allowbreak{}sqlite: cepas.\allowbreak{}cepa; desempate en la bitácora \\
\texttt{tipo\_vino} & atributo & — & — & Tinto o blanco. & origen\_vina.\allowbreak{}sqlite: cepas.\allowbreak{}tipo \\
\texttt{tier} & atributo & — & — & Varietal, Reserva o Gran Reserva. & CSV: tier \\
\texttt{formato} & atributo & — & — & 12 $\times$ 750 ml o 6 $\times$ 1500 ml. & CSV: formato \\
\texttt{litros\_por\_caja} & atributo & — & L por caja & Litros que rinde una caja: la definición de caja a litro. & CSV: litros\_por\_caja (9,0 para ambos formatos) \\
\texttt{anada} & atributo & — & año & Año de cosecha. & CSV: anada \\
\texttt{costo\_caja\_clp} & atributo & — & CLP por caja & Costo unitario de producción por caja. & CSV: costo\_caja\_clp; si falta, regla exacta (tier, añada) \\
\texttt{costo\_imputado} & atributo & — & 0/1 & 1 si el costo se imputó por regla. & generado \\
\texttt{activo} & atributo & — & 0/1 & Si el producto sigue en catálogo. & CSV: activo (si/SI/Sí/1 = verdadero) \\
\bottomrule
\end{longtable}
}

\subsubsection*{\texttt{dim\_cliente}}
Una fila por cliente (distribuidor o importador).

{\footnotesize
\begin{longtable}{@{}>{\raggedright\arraybackslash}p{0.18\textwidth}>{\raggedright\arraybackslash}p{0.09\textwidth}>{\raggedright\arraybackslash}p{0.08\textwidth}>{\raggedright\arraybackslash}p{0.075\textwidth}>{\raggedright\arraybackslash}p{0.2\textwidth}>{\raggedright\arraybackslash}p{0.245\textwidth}@{}}
\toprule
Columna & Rol & Aditiv. & Unidad & Descripción & Origen y regla \\
\midrule
\endhead
\texttt{cliente\_key} & llave & — & — & Llave del cliente. & generada \\
\texttt{cliente\_id} & atributo & — & — & Código del ERP. & origen\_vina.\allowbreak{}sqlite: clientes.\allowbreak{}cliente\_id \\
\texttt{nombre} & atributo & — & — & Razón social. & clientes.\allowbreak{}nombre \\
\texttt{tipo} & atributo & — & — & Distribuidor o importador. & clientes.\allowbreak{}tipo \\
\bottomrule
\end{longtable}
}

\subsubsection*{\texttt{dim\_mercado}}
Una fila por mercado de destino. País y zona son niveles de una misma jerarquía.

{\footnotesize
\begin{longtable}{@{}>{\raggedright\arraybackslash}p{0.18\textwidth}>{\raggedright\arraybackslash}p{0.09\textwidth}>{\raggedright\arraybackslash}p{0.08\textwidth}>{\raggedright\arraybackslash}p{0.075\textwidth}>{\raggedright\arraybackslash}p{0.2\textwidth}>{\raggedright\arraybackslash}p{0.245\textwidth}@{}}
\toprule
Columna & Rol & Aditiv. & Unidad & Descripción & Origen y regla \\
\midrule
\endhead
\texttt{mercado\_key} & llave & — & — & Llave del mercado. & generada \\
\texttt{pais} & atributo & — & — & País de destino. & origen\_vina.\allowbreak{}sqlite: mercados.\allowbreak{}pais \\
\texttt{zona} & atributo & — & — & Zona comercial. & mercados.\allowbreak{}zona\_id $\rightarrow$ zonas.\allowbreak{}zona; zona 8 inexistente (Corea del Sur) se asigna a Asia \\
\texttt{canal} & atributo & — & — & Nacional o Exportación. & mercados.\allowbreak{}canal \\
\texttt{hemisferio} & atributo & — & — & Norte o sur. & derivado del país (Chile y Brasil: sur) \\
\texttt{moneda\_facturacion} & atributo & — & — & CLP o USD. & monedas.\allowbreak{}moneda; en los datos la moneda queda determinada por el mercado (Chile CLP, resto USD); el ETL verifica que no haya excepciones \\
\bottomrule
\end{longtable}
}
